\chapter{Experimental Results}
\label{chap:results}
\setlength{\emergencystretch}{1.5em}

% ============================================================
\section{Experimental Overview}
\label{sec:results_overview}
% ============================================================

This chapter answers three questions: does RLInv provide measurable
operational advantages over existing rule-based and optimisation-based
approaches, under which delivery-logistics conditions do those
advantages emerge, and which design choices are responsible? It also
examines how far these conclusions hold under different reward, sizing,
and delivery-time assumptions. Evidence
is drawn from the full Phase 3 evaluation grid (10 sites, 10 seeds,
four delivery-logistics scenarios), reported primarily as Expected
Energy Not Served (EENS), with the site, scenario, and metric
definitions given in Chapter~\ref{chap:methodology}.

% ============================================================
\section{Primary Comparison}
\label{sec:primary_comparison}
% ============================================================

\emph{Does RLInv outperform existing approaches?}

Table~\ref{tab:primary_comparison} compares RLInv against four
reference policies. Two are rule-based: B0 uses fixed, uncalibrated
thresholds for dispatch and replenishment, whereas
B1 uses fixed default $(s,S)$ ordering thresholds
($s=0.25$, $S=0.85$) with rule-based dispatch. One is optimisation-based: a rolling-horizon
Model Predictive Control baseline (MPC) that addresses the standing
question of whether classical optimisation is sufficient once its
required forecasts are supplied. The fourth (TrackB) pairs an analytically
parameterised $(s,S)$ ordering rule with a
\emph{learned} (reinforcement-learning)
dispatch policy, and is included to examine a distinct question from
the other three: whether reinforcement learning \emph{alone}, applied
only to dispatch, is sufficient, or whether the reliability gains
reported below require learning ordering and dispatch \emph{jointly}.
An oracle variant of MPC, given perfect foresight of load, solar
generation, and grid availability, is included as an optimistic reference representing the
best achievable dispatch performance within the decomposition
framework under perfect future information (diesel delivery timing remains
stochastic and governed by the same classical ordering rule as the
other MPC variants); it is not intended as a
deployable policy and is not counted among the competing baselines.

\begin{table}[htb]
\centering
\small
\caption{Primary comparison: mean EENS (kWh) by policy and delivery
scenario, aggregated over 10 sites $\times$ 10 seeds. Lower is better.}
\label{tab:primary_comparison}
\begin{tabular}{p{5.2cm}rrrr}
\toprule
Policy & Normal & Delayed & Monsoon & Extreme \\
\midrule
B0 (threshold rule) & 9.94 & 24.53 & 35.52 & 66.63 \\
B1 ($(s,S)$ + rule dispatch) & 3.55 & 13.43 & 23.20 & 57.73 \\
MPC ($H{=}24$, persistence) & 3.08 & 13.14 & 22.84 & 57.38 \\
TrackB ($(s,S)$ + learned dispatch) & 1.28 & 3.65 & 174.50 & 152.13 \\
\midrule
\textit{Oracle-MPC (perfect forecast, reference only)} & \textit{3.17} & \textit{13.26} & \textit{23.20} & \textit{57.45} \\
\midrule
\textbf{RLInv (proposed, joint learning)} & \textbf{0.02} & \textbf{3.60} & \textbf{11.29} & \textbf{53.85} \\
\bottomrule
\end{tabular}
\end{table}

\revised{Within the original-configuration policy set in
Table~\ref{tab:primary_comparison}, RLInv records the lowest
mean EENS under all four delivery-logistics scenarios.
That table excludes the separately calibrated B1*
variants in Subsection~\ref{sec:calibrated_baselines};
B1*-HighBuffer records lower numerical mean EENS
than RLInv in three stressed scenarios.} Under normal
logistics, most policies perform adequately, but the advantage of
RLInv widens substantially as delivery uncertainty increases: relative to the original fixed-threshold B1 baseline, RLInv reduces EENS by 73\%
under delayed logistics, 51\% under monsoon conditions, and 7\% under
extreme logistics stress. The narrowing margin under the most severe
scenario is itself informative: under a mean 336-hour delivery lead
time, the achievable planning horizon for \emph{any} policy is bounded
by physical fuel availability, and the residual advantage reflects
RLInv's more efficient scheduling of a scarce resource rather than a
qualitatively different capability (see Section~\ref{sec:mechanism_h1}).

The comparison with Oracle-MPC, on paired (site, seed) comparisons
($n{=}100$), indicates a real difference under normal, delayed, and
monsoon conditions: mean paired differences of 3.15, 9.67, and
11.91\,kWh respectively, with 95\% confidence intervals excluding zero
in all three cases. Under extreme conditions the paired difference is
directionally consistent but the interval spans zero (3.60\,kWh, 95\%
CI $[-0.18, 7.38]$). This gap is largest under
Monsoon but narrows again under Extreme, where physical fuel
availability increasingly constrains all policies --- indicating that RLInv's
benefit cannot be explained solely by improved forecasting or dispatch
optimisation, since Oracle-MPC already has access to realised future
load, solar generation, and grid availability for dispatch planning.

This reliability advantage is not obtained through disproportionate
diesel consumption: aggregated across all sites and scenarios, RLInv
consumes 205.7\,kWh of diesel versus B1's 198.5\,kWh (approximately
3.6\% more), a modest increase consistent with a genuine operational
trade-off rather than an inflated or unfair comparison.
Table~\ref{tab:secondary_metrics} shows that this modest increase is
accompanied by a consistently larger inventory buffer and more proactive
replenishment behaviour.

\begin{table}[htb]
\centering
\small
\caption{Secondary operational metrics by policy and scenario, aggregated
over 10 sites $\times$ 10 seeds.}
\label{tab:secondary_metrics}
\begin{tabular}{llrrrrr}
\toprule
\textbf{Policy} & \textbf{Scenario} & \textbf{EENS} & \textbf{Diesel} &
\textbf{Completed} & \textbf{Mean} & \textbf{Minimum} \\
 & & (kWh) & (kWh) & \textbf{orders/episode} & \textbf{inventory (\%)} & \textbf{inventory (\%)} \\
\midrule
RLInv  & Normal  & 0.02   & 222.9 & 2.16 & 78.8 & 52.4 \\
RLInv  & Monsoon & 11.29  & 211.6 & 1.97 & 68.1 & 45.4 \\
RLInv  & Extreme & 53.85  & 169.1 & 1.34 & 53.1 & 37.3 \\
\midrule
B1     & Normal  & 3.55   & 219.4 & 0.63 & 54.0 & 36.6 \\
B1     & Monsoon & 23.20  & 199.7 & 0.58 & 50.9 & 35.2 \\
B1     & Extreme & 57.73  & 165.2 & 0.50 & 47.0 & 34.4 \\
\midrule
TrackB & Normal  & 1.28   & 221.7 & 1.13 & 53.2 & 39.0 \\
TrackB & Monsoon & 174.50 & 48.4  & 0.90 & 74.3 & 55.8 \\
TrackB & Extreme & 152.13 & 70.8  & 0.79 & 61.8 & 50.6 \\
\midrule
A6     & Normal  & 1.28   & 221.7 & 1.13 & 53.2 & 39.0 \\
A6     & Monsoon & 10.65  & 212.3 & 3.79 & 75.0 & 47.1 \\
A6     & Extreme & 61.50  & 161.4 & 0.69 & 46.2 & 34.6 \\
\bottomrule
\end{tabular}
\end{table}

RLInv's ordering frequency and inventory buffer both decline
monotonically as logistics stress increases (orders: 2.16 $\to$ 1.34;
mean inventory: 78.8\% $\to$ 53.1\%, Normal to Extreme), consistent with
proactive replenishment rather than reactive over-ordering under stress.
TrackB's low diesel consumption under Monsoon and Extreme should not be
interpreted as greater efficiency; instead, poor dispatch leaves more
fuel unused while EENS increases substantially. Under Normal logistics,
TrackB and A6 produce nearly identical outcomes because both share the
same calibrated $(s,S)$ replenishment policy and dispatch decisions have
little influence under low-stress conditions.

Notably, the TrackB comparison shows that replacing rule-based dispatch with \emph{learned} dispatch
while keeping classical $(s,S)$ ordering does not
reproduce RLInv's advantage --- indeed it performs \emph{worse} than
either rule-based baseline under monsoon and extreme conditions. This
suggests that reinforcement learning alone, applied
only to the dispatch decision, does not explain RLInv's performance.

MPC approximately optimises dispatch over a 24-hour receding horizon using the available forecast, yet still under-performs
RLInv because its planning horizon ($H{=}24$\,h) cannot see far enough
ahead to plan for delivery delays that can extend far beyond the
24-hour planning horizon.

\textbf{Improving dispatch alone --- whether through reinforcement
learning or through optimisation --- does not eliminate the degradation
caused by externally managed replenishment.} RLInv achieves the best
overall performance in the evaluated grid (Figure~\ref{fig:eens_comparison}); the following ablations
examine how much of that result can be attributed to joint replenishment
learning and the other design choices. This motivates the mechanistic analysis in
Section~\ref{sec:mechanism}, where the three pre-specified ablation
studies isolate the individual contributions of coupled replenishment
learning, action masking, and inventory-aware state representation.

\begin{figure}[htb]
\centering
\includegraphics[width=\textwidth]{fig_eens_comparison_v4.png}
\caption{Overall policy performance across delivery scenarios. Policies,
left to right: \textbf{B0} (threshold dispatch, fixed uncalibrated
threshold ordering), \textbf{B1} (threshold dispatch, ixed-threshold $(s,S)$
ordering), \textbf{MPC} (model-predictive dispatch; planning horizon
discussed in Chapter~\ref{chap:methodology}), \textbf{TrackB}
($(s,S)$ ordering with a separately learned dispatch policy),
\textbf{Oracle-MPC} (dispatch with perfect future load, solar and grid
information; an optimistic reference representing the best achievable
dispatch performance within the decomposition framework, not a
deployable policy; diesel delivery remains stochastic), and \textbf{RLInv} (the proposed joint
policy). Error bars show 95\% confidence intervals over $n=100$
(10 sites $\times$ 10 seeds).}
\label{fig:eens_comparison}
\end{figure}

\subsection{Performance across sites with different operating characteristics}
\label{sec:site_breakdown}

The comparison above aggregates over all ten sites; it is worth
checking that RLInv's advantage is not driven by a small number of
sites while performing poorly elsewhere. Figure~\ref{fig:site_breakdown}
reports RLInv and B1 mean EENS under normal logistics, computed
separately for each site (each site's policy is trained and evaluated
on that site individually), grouped by an independent difficulty
classification (solar coverage, outage frequency, and battery
autonomy; see Chapter~\ref{chap:datasets-setup},
Section~\ref{sec:site_classification} for the classification
criteria). RLInv matches or improves on
B1 at every site, among the largest absolute gains on the hardest
(lowest-autonomy, highest-outage) sites, where B1's reactive ordering
is least able to keep pace with demand. \begin{revisedblock}
This is a per-site robustness check, not a test
of generalisation to unseen sites: each policy
in this comparison was trained for its own site.
Cross-site transfer without site-specific
retraining is evaluated separately in
Subsection~\ref{sec:cross_site_transfer}.
\end{revisedblock}

\begin{figure}[htb]
\centering
\includegraphics[width=\textwidth]{fig_cross_site.png}
\caption{Performance across sites with different operating
characteristics: mean EENS (kWh/episode) for RLInv and B1, each
trained and evaluated on its own site, normal logistics. X-axis label
colour indicates the independent difficulty classification (green =
Easy, orange = Medium, red = Hard).}
\label{fig:site_breakdown}
\end{figure}

% ============================================================
\section{Understanding Why RLInv Works}
\label{sec:mechanism}
% ============================================================

The primary comparison establishes that RLInv consistently achieves
the best overall performance. To understand which design choices
contribute to this, three planned ablation comparisons examine the
contributions of replenishment learning, action masking, and
logistics-related observation channels (Chapter~\ref{chap:methodology}).
Not every comparison changes only one implementation detail, so the
limitations of each comparison are stated explicitly below. They are referred to here by what they test --- \textbf{joint
replenishment and dispatch learning}, \textbf{action masking}, and
\textbf{inventory observability} --- corresponding to the H1, H2, and
H3 labels used in the methodology chapter and analysis scripts. All
comparisons below are
\emph{paired at the seed level}: the same 10 seeds are evaluated under
every policy, so seed 42 under RLInv is matched against seed 42 under
the comparator, not treated as an independent sample. Significance is
therefore assessed with a paired $t$-test on the within-seed
differences, with effect size Cohen's $d_z$ (mean paired difference /
std of paired differences), following the corrected statistical
methodology adopted after Phase 2 (an earlier, independent-samples test
understated significance for this design).

\subsection{Joint replenishment and dispatch learning}
\label{sec:mechanism_h1}

\textbf{Evidence.} A6 replaces jointly learned replenishment with an
external calibrated $(s,S)$ policy and re-trains dispatch using a
reduced two-action space (\texttt{Discrete(2)}, generator on/off only)
while retaining the full 11-dimensional observation, in contrast to
RLInv's six-action joint space (\texttt{Discrete(6)}) (Figure~\ref{fig:ablation_h1}). It therefore
provides a cleaner replenishment comparison than TrackB, but does not
hold the dispatch architecture fully constant --- it does not isolate
``ordering'' as a standalone mechanism, since A6
still learns its own dispatch policy, under a different action space
than RLInv's joint action. The paired comparison is scenario-dependent rather than
uniform: RLInv is significantly better under Normal ($p=0.0175$,
$d_z=-0.92$) and Extreme ($p=0.00035$, $d_z=-1.76$, a large effect),
the difference is negligible under Delayed ($p=0.83$, $d_z=-0.07$),
and under Monsoon the comparison is not statistically distinguishable
and nominally favours A6 ($p=0.41$, $d_z=+0.27$). The same qualitative
pattern holds when restricted to the selected three-site sensitivity
subset (site2, site5, and site7), so
this is not an artefact of averaging across easier sites.

\textbf{Interpretation.} The additional benefit attributable to
jointly learned replenishment was observed under Normal and Extreme
logistics, but not under Delayed or Monsoon conditions. Diagnostic
inspection (constraint violations, stockout events,
diesel usage; Section~\ref{sec:behaviour}) shows zero hard violations
for both policies throughout, so the Normal/Extreme difference is not
explained by infeasible or erratic actions; it is consistent with
RLInv's ability to anticipate that its own dispatch decisions can
deplete inventory ahead of a delayed delivery, a benefit that is not
uniformly present across every logistics regime tested.

For context, RLInv also differs from TrackB (classical $(s,S)$
ordering paired with a \emph{separately trained} dispatch policy, a
more confounded comparator than A6 since both ordering and dispatch
architecture differ). This gap is large and consistent
(Table~\ref{tab:primary_comparison}: e.g.\ 11.29 vs.\ 174.50\,kWh
monsoon, 53.85 vs.\ 152.13\,kWh extreme), but because TrackB confounds
two design differences at once, it is used below to examine
masking-related architectural differences, with the observation-space
confound stated explicitly, rather than as primary evidence for joint
replenishment.

\begin{figure}[htb]
\centering
\includegraphics[width=0.98\textwidth]{fig_ablation_h1_v3.png}
\caption{Joint replenishment and dispatch learning: RLInv vs.\ A6 (cleaner
practical replenishment comparison; dispatch is re-trained under a
reduced action space) and vs.\ TrackB
(shown for context; confounds ordering and dispatch architecture). Paired
$t$-test significance shown where applicable: * $p<0.05$, ***
$p<0.001$, ns = not significant. Error bars: 95\% CI, $n=10$ seeds.}
\label{fig:ablation_h1}
\end{figure}

\subsection{Action masking}
\label{sec:mechanism_h2}

\textbf{Evidence.} A6 (masked dispatch) and TrackB (unmasked
dispatch) share an identical classical $(s,S)$ ordering rule; A6 retains
the full 11-dimensional observation while TrackB's dispatch agent uses a
reduced 6-dimensional observation (Chapter~\ref{chap:methodology}), so
this comparison differs in both masking and observation design, not
masking alone. There is no meaningful difference
under Normal or Delayed logistics (identical or not significant), but a
very large, highly significant effect under Monsoon ($d_z=-10.2$,
$p<0.0001$) and Extreme ($d_z=-4.6$, $p<0.0001$), both favouring the
masked (A6) variant (Figure~\ref{fig:ablation_h2}). Separately, RLInv vs.\ A7 (vanilla PPO, no
masking, otherwise identical to RLInv, so ordering is already jointly
learned in both, and both share the same 11-dimensional observation)
shows no significant difference in any scenario
($|d_z|\leq 0.37$ throughout).

\textbf{Interpretation.} The A6--TrackB result is consistent with an
architecture-dependent robustness difference under stressed logistics,
but cannot be attributed to masking alone because the observation
designs also differ; the resulting gap is also confounded with
whether ordering is jointly learned, since RLInv and TrackB differ in
that respect too. The defensible claim is that masking is associated
with a substantial part of the RLInv--TrackB gap under Monsoon/Extreme,
not all of it. The cleaner RLInv-vs-A7 comparison (identical
observation space, ordering already jointly learned in both) shows no
statistically distinguishable mean-EENS effect from masking in any
scenario. Both A6 and TrackB record zero explicit hard constraint
violations, so the A6--TrackB effect is not about constraint
enforcement \emph{per se}. Masking's benefit is therefore conditional
on ordering being externally managed, not a general main effect of the
learning algorithm.

\begin{figure}[htb]
\centering
\includegraphics[width=0.85\textwidth]{fig_ablation_h2_v3.png}
\caption{Action masking: A6 (masked) vs.\ TrackB (unmasked); the external
$(s,S)$
ordering policy is identical in both (observation design differs --- see
text). Effect is large and significant under
Monsoon/Extreme only. See text for why this is not interpreted as
explaining the entire RLInv--TrackB gap, and for the separate,
non-significant RLInv-vs-A7 comparison.}
\label{fig:ablation_h2}
\end{figure}

\subsection{Inventory observability}
\label{sec:mechanism_h3}

\textbf{Evidence.} A5 is trained identically to RLInv but with five
logistics/inventory-related channels removed from its observation
space (inventory level, pending-delivery flag, pending quantity,
elapsed time since ordering, and the remaining delivery-time channel).
RLInv achieves numerically lower mean EENS than A5 in every
scenario, but the paired differences are small and do not reach
statistical significance under the final evaluation protocol
($|d_z|$ ranging 0.24--0.69; closest to significance under Extreme,
$p=0.058$, still not significant at $\alpha=0.05$) (Figure~\ref{fig:ablation_h3}).

\textbf{Interpretation.} Inventory observability produced consistently
lower mean EENS, but the differences remained small and were not
statistically distinguishable under the final evaluation protocol. This
is reported as a genuine finding rather than omitted, since it
clarifies which state information the learned policy actually exploits
--- the agent appears able to infer much of the required information
indirectly from other observable variables (battery state-of-charge,
grid availability, load, solar generation, and time-of-day information).
The combined contribution of the removed logistics-related channels
appears modest under the final evaluation protocol.

\begin{figure}[htb]
\centering
\includegraphics[width=0.75\textwidth]{fig_ablation_h3_v3.png}
\caption{Inventory observability: RLInv (observes tank fill fraction,
pending-delivery status, and ordering-history channels) vs.\ A5
(identical, these five logistics/inventory-related observations
zeroed). RLInv is numerically lower in every scenario but the
difference does not reach significance ($p=0.058$ at closest, Extreme
scenario). Error bars: 95\% CI, $n=10$ seeds.}
\label{fig:ablation_h3}
\end{figure}



% ============================================================
\section{Behavioural Analysis}
\label{sec:behaviour}
% ============================================================

A representative single-site trajectory (Figure~\ref{fig:trajectory})
illustrates patterns also visible in aggregate across the full
evaluation grid. Mean ordering frequency (RLInv, count of orders
actually placed after action masking and tank-capacity clipping ---
not the number of ordering actions attempted) decreased monotonically
from 2.16 orders per episode under Normal logistics to 1.34 under
Extreme disruption. This is consistent with the policy reducing new
orders when delivery completion within the remaining evaluation
horizon becomes increasingly unlikely; 
diesel generator use is
curtailed selectively rather than run to exhaustion.

\revised{Across the 44,000 archived evaluation episodes, no
stockout events or physical-clipping violations were recorded,
although unmet load occurred. The corresponding stockout and
violation penalties were therefore inactive during evaluation.}


These behavioural observations provide an explanatory interpretation of
the quantitative differences identified in the preceding hypothesis
analyses rather than constituting additional independent statistical
evidence.

\begin{figure}[htb]
\centering
\includegraphics[width=\textwidth]{fig_trajectory.png}
\caption{Representative single-episode trajectory (site2, monsoon
scenario): inventory fraction, diesel orders placed, and unmet energy
demand over the episode, compared between RLInv and B1.}
\label{fig:trajectory}
\end{figure}

% ============================================================
\section{Additional Experimental Evaluation}
\label{sec:robustness}
% ============================================================

This section tests whether the preceding conclusions survive changes
to reward design, physical sizing, and delivery-time modelling.

\subsection{Sensitivity to delivery lead time}
\label{sec:lead_sensitivity}

The four scenarios used throughout this chapter (Normal, Delayed,
Monsoon, Extreme) are discrete operating points; this subsection
instead asks how performance evolves across a wider sequence of mean
delivery lead-time settings, increasing from 0h to 336h, \revised{evaluated across sites 2, 5, 7, and 10
using the evaluation seeds} (Figure~\ref{fig:lead_sensitivity}).

\begin{figure}[htb]
\centering
\includegraphics[width=\textwidth]{fig_lead_sensitivity.png}
\caption{Policy performance under increasing delivery delays. (a) Mean
EENS vs.\ mean delivery lead time for RLInv, B1, and B0: all three
degrade as delays increase, with RLInv consistently below B1 and B0
consistently worst. (b) RLInv's absolute advantage over B1 across the
same range. The advantage grows through 72h, then narrows from 120h
onward --- not because RLInv stops working, but because at very long
lead times physical tank capacity becomes the binding constraint for
every policy: there is only so much fuel that can be pre-ordered
regardless of ordering strategy, so the RLInv--B1 gap narrows as
physical tank capacity becomes increasingly binding.}
\label{fig:lead_sensitivity}
\end{figure}

RLInv's advantage over B1 is small at short lead times (1.3\,kWh at
24h), peaks at 29.5\,kWh around 72h, and narrows to 8.6\,kWh by 336h
even though RLInv remains numerically better throughout (134.3 vs.\
142.9\,kWh at 336h). This wider sequence of lead-time settings is consistent with, and adds
resolution to, the four discrete scenarios reported elsewhere in this
chapter.

\subsection{Reward Function Sensitivity}
Ten reward-weight configurations were evaluated, independently varying
the unmet-demand penalty ($\lambda \in \{25,\dots,500\}$, a 20-fold
range), the stockout penalty, the grid-cost weight, and the constraint
violation penalty. Extreme-scenario EENS varied by only 3.32\,kWh
(167.74--171.06\,kWh) across the entire sweep, and Monsoon-scenario
EENS varied by 2.90\,kWh across the sweep. Separately, training
with $\gamma \in \{0.95, 0.99, 0.995\}$ produced results within
2.4\,kWh of one another on every scenario (Figure~\ref{fig:reward_sensitivity}).

\begin{figure}[htb]
\centering
\includegraphics[width=\textwidth]{fig_reward_sensitivity.png}
\caption{Sensitivity of RLInv's performance to reward weights (a) and
discount factor
(b), evaluated on the selected
three-site sensitivity
subset under the monsoon and extreme scenarios, 3 seeds $\times$ 400k training
steps.}
\label{fig:reward_sensitivity}
\end{figure}

RLInv's EENS remained within a narrow range across the tested
reward-weight and discount-factor settings, indicating that its own
performance is not highly sensitive to these choices.

\subsection{Tank capacity sensitivity}
RLInv's advantage over B1 was evaluated across five tank capacities
(24\,h, 48\,h, 72\,h baseline, 144\,h, and 336\,h of buffer) on the
selected three-site sensitivity subset under delayed logistics
(Table~\ref{tab:tank_sensitivity}, Figure~\ref{fig:tank_sensitivity}).

\begin{table}[htb]
\centering
\caption{Tank capacity sensitivity: mean EENS (kWh), selected
three-site sensitivity subset,
delayed scenario, $n=10$ seeds.}
\label{tab:tank_sensitivity}
\begin{tabular}{lrrr}
\toprule
Tank capacity & RLInv & B1 & RLInv advantage \\
\midrule
24\,h (0.33$\times$) & 215.554 & 283.553 & 67.999 \\
48\,h (0.67$\times$) & 42.876 & 94.682 & 51.806 \\
72\,h (1.0$\times$, baseline) & 11.949 & 44.630 & 32.681 \\
144\,h (2.0$\times$) & 0.000 & 4.812 & 4.812 \\
336\,h (4.67$\times$) & 0.000 & 0.000 & 0.000 \\
\bottomrule
\end{tabular}
\end{table}

\begin{figure}[htb]
\centering
\includegraphics[width=0.8\textwidth]{fig_tank_sensitivity.png}
\caption{RLInv's advantage over B1 as a function of tank capacity,
relative to each scenario's mean delivery lead time. The advantage
vanishes once tank capacity is large enough to absorb delivery
uncertainty without proactive ordering.}
\label{fig:tank_sensitivity}
\end{figure}

RLInv's advantage decreases monotonically as tank capacity grows, from
68.0\,kWh at 24\,h to 0\,kWh at 336\,h: inventory coupling matters
specifically when tank capacity is small relative to delivery lead
time, and a reactive rule-based policy becomes adequate once the
physical buffer is large enough to absorb the same uncertainty.

\subsection{Distributional robustness}
The evaluation grid above assumes delivery lead times follow a
geometric distribution. To test sensitivity to this modelling choice,
additional evaluation runs were conducted under three alternative
delivery-time distributions with matched scenario means: lognormal
($\sigma=0.5$), lognormal ($\sigma=0.8$, heavier tail), and Weibull
($k=2$) (Figure~\ref{fig:distributional_robustness}).

\begin{figure}[htb]
\centering
\includegraphics[width=0.85\textwidth]{fig_distributional_robustness.png}
\caption{EENS under the training (geometric) delivery-time
distribution and three alternative distributions --- lognormal
($\sigma=0.5$), lognormal ($\sigma=0.8$), and Weibull ($k=2$) --- for
RLInv, B1, and MPC across all four scenarios. Log-scaled $y$-axis to
accommodate the extreme-scenario range.}
\label{fig:distributional_robustness}
\end{figure}

Policy rankings are preserved across all four scenarios and all three
alternative distributions, and the mean absolute change in EENS
relative to the training (geometric) distribution is smaller for RLInv
than for B1 in the lognormal comparisons (4.56\,kWh vs.\ 7.62\,kWh at
$\sigma=0.5$).

The preserved rankings indicate that the main comparison is not
sensitive to the particular lead-time distribution tested, and if
anything RLInv is less sensitive to this modelling choice than the
classical baseline.

\subsection{Stochastic grid outage}
This comparison evaluates RLInv under a stochastic (Markov-chain) grid
outage model rather than the deterministic weekly template used
elsewhere in this chapter. Both use the same sites, training budget,
seeds and evaluation protocol, but differ in the grid process used
during training. The two policies are evaluated on the same three-site
sensitivity subset
(site2, site5, site7), the same 3 seed realisations, and the same 10
episodes per (site, scenario) cell in both cases --- a smaller
evaluation grid than the 10-site comparison used elsewhere in this
chapter, since this experiment was run only on the sites where
logistics stress is most likely to be informative.

\begin{table}[htb]
\centering
\caption{EENS (kWh) under the stochastic grid outage model:
deterministic-grid-trained policy evaluated on the stochastic grid,
vs.\ a policy trained and evaluated on the stochastic grid. 3 hard
sites, 3 seeds, 10 episodes per cell.}
\label{tab:stochastic_grid}
\begin{tabular}{lrrr}
\toprule
Scenario & Deterministic-trained (eval on stochastic) & Stochastic-trained & Change \\
\midrule
Normal & 2.96 & 0.24 & $-$91.9\% \\
Delayed & 21.49 & 27.30 & $+$27.0\% \\
Monsoon & 42.62 & 61.12 & $+$43.4\% \\
Extreme & 153.27 & 181.08 & $+$18.1\% \\
\bottomrule
\end{tabular}
\end{table}

The policy trained under stochastic grid outages achieved lower EENS
under Normal conditions but higher EENS under Delayed, Monsoon and
Extreme logistics. The largest relative degradation
occurs under the Monsoon scenario, indicating a trade-off between
adapting to grid uncertainty and conserving diesel under prolonged
delivery disruption. This is reported as an identified limitation
rather than resolved, and, given the narrower evaluation grid used
here, is not generalised beyond the three sites tested.

\subsection{RLInv-DA: exploratory extension motivated by reviewer feedback}
\label{sec:rlinv_da}

Reviewer feedback received during the Phase~2 evaluation raised two
related questions: whether the diesel supply chain could be modelled
more realistically, including vendor-side considerations, and whether
the reordering strategy could track and
predict delivery delays rather than treating each delivery as an
independent random event. This subsection reports an
exploratory extension, RLInv-DA, introduced to investigate delay
awareness within the stated single-vendor abstraction, without altering the main
results reported in Sections~\ref{sec:primary_comparison}--\ref{sec:mechanism}.

RLInv-DA extends RLInv with a single additional observation,
$\widehat{L}_t$, an exponentially weighted moving average of realised
delivery delays updated on each arrival ($\widehat{L}_t = 0.3\,
L_t^{\mathrm{realised}} + 0.7\, \widehat{L}_{t-1}$), evaluated under a
two-state (Normal/Disrupted) supplier-regime environment that gives
delivery delays the temporal correlation an EWMA needs to have
anything to detect. Alternative approaches, including predictive
lead-time models and more detailed supply-chain formulations, were
considered during the design phase but were not adopted because they
addressed different problem formulations or required assumptions
beyond the scope of this study (see Chapter~\ref{chap:methodology}).
Both RLInv-DA and a base-regime RLInv policy were
retrained from scratch under identical conditions, and diagnostic
trace analysis confirmed the estimator updates correctly and that
ordering decisions consistently occur after each update --- the policy
has genuine opportunity to use the signal.

\textbf{Results.} Unless otherwise stated, results in this subsection
are aggregated over 3 sites $\times$ 3 training seeds $\times$ 4
scenarios $\times$ 10 evaluation episodes per variant (720 episodes
total). Table~\ref{tab:rlinv_da_results} and Figure~\ref{fig:rlinv_da} report EENS across
all four scenarios for both variants. Confidence intervals overlap in every scenario, and no
consistent direction of improvement emerges: across the full grid of
36 site--seed--scenario comparisons, RLInv-DA achieved lower EENS in 5
cases, higher EENS in 11 cases, and statistically indistinguishable
performance in the remaining 20. Restricting to the two scenarios
where delivery-delay tracking is most plausibly informative (delayed
and monsoon), the pattern is unchanged (3 better, 5 worse, 10
comparable of 18), and an extended evaluation across normal and
extreme logistics confirms this is not an artefact of scenario choice
rather than a genuine absence of consistent benefit: both policies are
indistinguishable under normal logistics, where the control problem is
too easy to differentiate, and differ by under 1\% under extreme
logistics, where both policies incur high EENS regardless of delay
awareness.

\begin{table}[htb]
\centering
\caption{RLInv-DA performance: mean EENS (kWh) with 95\% confidence
intervals, aggregated over 3 sites $\times$ 3 seeds.}
\label{tab:rlinv_da_results}
\begin{tabular}{lrr}
\toprule
Scenario & Base-regime RLInv & RLInv-DA (EWMA) \\
\midrule
Normal & 0.000 & 0.000 \\
Delayed & 4.408 [0.86, 7.95] & 3.779 [0.13, 7.43] \\
Monsoon & 20.297 [10.22, 30.37] & 23.111 [12.37, 33.86] \\
Extreme & 233.632 [202.2, 265.0] & 235.429 [204.0, 266.9] \\
\bottomrule
\end{tabular}
\end{table}

\begin{figure}[htb]
\centering
\includegraphics[width=\textwidth]{fig_rlinv_da.png}
\caption{RLInv-DA vs.\ base-regime RLInv: mean EENS (kWh) with 95\%
confidence intervals across all four scenarios, 3 sites $\times$ 3
seeds.}
\label{fig:rlinv_da}
\end{figure}

Behavioural analysis across 72 matched representative trace episodes
(3 sites $\times$ 3 seeds $\times$ 4 scenarios $\times$ 2 variants,
one representative episode per site--seed--scenario combination rather
than an average over the full 10-episode evaluation grid used for the
performance results above)
(Table~\ref{tab:rlinv_da_behaviour}, Figure~\ref{fig:rlinv_da_behaviour}) shows that the EWMA-aware policy
is behaviourally active rather than ignoring the additional
observation, but its effect is scenario-dependent and not uniformly
protective. Under monsoon conditions, RLInv-DA orders later and at
lower inventory than the base-regime policy (mean inter-order interval
150.0\,h versus 137.2\,h, +9.4\%); under extreme conditions the effect
reverses, with RLInv-DA ordering slightly earlier and at higher
inventory (160.0\,h versus 164.3\,h). This reversal is consistent with
an estimator that responds to the overall level of realised delay
specific to each regime: the persistently long delays characteristic
of extreme logistics push the EWMA estimate upward, prompting earlier
ordering, whereas the comparatively shorter and more variable delays
under monsoon conditions do not sustain the same elevated estimate,
occasionally licensing later ordering that increases exposure.

\begin{table}[htb]
\centering
\caption{RLInv-DA behavioural mechanism: mean inventory fraction at
order placement and mean inter-order interval, 72 matched representative
trace episodes (one per site--seed--scenario combination, not a
10-episode average).}
\label{tab:rlinv_da_behaviour}
\begin{tabular}{lrrrr}
\toprule
& \multicolumn{2}{c}{Inventory at order} & \multicolumn{2}{c}{Inter-order interval (h)} \\
Scenario & Base & EWMA & Base & EWMA \\
\midrule
Normal & 0.650 & 0.669 & 101.5 & 103.4 \\
Delayed & 0.628 & 0.628 & 108.3 & 112.6 \\
Monsoon & 0.588 & 0.566 & 137.2 & 150.0 \\
Extreme & 0.596 & 0.603 & 164.3 & 160.0 \\
\bottomrule
\end{tabular}
\end{table}

\begin{figure}[htb]
\centering
\includegraphics[width=\textwidth]{fig_rlinv_da_behaviour.png}
\caption{RLInv-DA behavioural mechanism: inventory fraction at order
placement and inter-order interval, base-regime vs.\ EWMA-aware
policy, across the four scenarios (72 matched representative trace
episodes).}
\label{fig:rlinv_da_behaviour}
\end{figure}

\textbf{Interpretation.} The contribution of this extension is not
that explicit delay estimation improved reliability. It is that,
under a controlled supplier-regime model designed specifically to
give such estimation a genuine opportunity to help, the mechanism was
implemented, verified, and evaluated, and found to provide no
consistent additional benefit beyond the information already available to RLInv
through its existing inventory and pipeline state. This is reported as
a completed finding addressing the underlying reviewer feedback,
rather than deferred as future work.

% rlinv_da_design_rationale.tex
% Research workflow diagram for RLInv-DA exploratory extension (Comments #9/#10)
% REDESIGNED: two-column zigzag layout (was single tall column, forced into
% resizebox{!}{0.78\textheight} which crushed font size and box width).
% This version scales to full page width instead, at a near-square aspect
% ratio, so boxes and text render at a comfortable size.
%
% Required packages (add to thesis preamble if not already present):
%   \usepackage{tikz}
%   \usetikzlibrary{arrows.meta, shapes.geometric, positioning, fit, backgrounds}
%   \usepackage{xcolor}

\begin{figure}[htbp]
\centering
\resizebox{\textwidth}{!}{%
\begin{tikzpicture}[
    node distance = 7mm and 14mm,
    % Box styles
    reviewer/.style = {
        rectangle, rounded corners=3pt,
        draw=black!60, fill=gray!10,
        text width=68mm, align=center,
        font=\normalsize\itshape, inner sep=5pt
    },
    problem/.style = {
        rectangle, rounded corners=3pt,
        draw=orange!70, fill=orange!8,
        text width=68mm, align=center,
        font=\normalsize, inner sep=5pt
    },
    decision/.style = {
        rectangle, rounded corners=3pt,
        draw=blue!50, fill=blue!6,
        text width=68mm, align=center,
        font=\normalsize, inner sep=5pt
    },
    result/.style = {
        rectangle, rounded corners=3pt,
        draw=teal!60, fill=teal!6,
        text width=68mm, align=center,
        font=\normalsize, inner sep=5pt
    },
    conclusion/.style = {
        rectangle, rounded corners=3pt,
        draw=green!50!black, fill=green!6,
        text width=140mm, align=center,
        font=\normalsize\bfseries, inner sep=6pt
    },
    evidence/.style = {
        rectangle, rounded corners=3pt,
        draw=purple!50, fill=purple!5,
        text width=68mm, align=left,
        font=\normalsize, inner sep=6pt
    },
    % Arrow style
    arr/.style = {-{Stealth[length=6pt]}, thick, gray!70},
    jump/.style = {-{Stealth[length=6pt]}, thick, gray!50, dashed}
]

% Column 1: reviewer comment through motivation
\node[reviewer] (comment)
    {Research Question: Delay-Aware Replenishment\\[2pt]
     \textit{Can past delivery delays improve
     diesel replenishment decisions?}};

\node[problem, below=of comment] (iid)
    {Original environment: i.i.d.\ geometric delivery times\\[2pt]
     \textrightarrow\ Past delays carry no predictive information};

\node[decision, below=of iid] (regime)
    {\textbf{Supplier model:} Two-state Markov regime\\[2pt]
     Normal $\leftrightarrow$ Disrupted (persistent disruptions)\\
     Disrupted mean $= 2\times$ normal mean};

\node[problem, below=of regime] (informative)
    {Past realised delays now contain predictive information\\[2pt]
     \textrightarrow\ EWMA estimation is scientifically motivated};

% Column 2: response through evaluation
\node[decision, right=of comment, xshift=14mm] (ewma)
    {\textbf{Delay-history observation:} EWMA estimator\\[2pt]
     $\hat{L}_{t} = \alpha \cdot L^{\text{realised}}_t + (1-\alpha)\cdot\hat{L}_{t-1}$, \quad $\alpha = 0.3$\\
     No future leakage; uses only past observations};

\node[decision, below=of ewma] (rlinvda)
    {\textbf{RLInv-DA:} Base observation space $+$ EWMA estimate\\[2pt]
     Both \textit{Base-regime} and \textit{EWMA-aware} retrained\\
     under same supplier model (fair ablation)};

\node[result, below=of rlinvda] (eval)
    {\textbf{Evaluation}\\[2pt]
     Performance: 720 episodes (3 sites $\times$ 3 seeds $\times$ 4 scenarios)\\
     Behaviour: 72 matched trace episodes\\
     Diagnostics: formula verified, regime confirmed, signal available};

\node[evidence, below=of eval] (evidence)
    {\textbf{Evidence}\\[3pt]
     {\color{green!50!black}$\checkmark$} EWMA formula correct ($\alpha = 0.3$, all updates verified)\\
     {\color{green!50!black}$\checkmark$} Supplier regime active (transitions confirmed)\\
     {\color{green!50!black}$\checkmark$} Bidirectional adaptation (up and down)\\
     {\color{green!50!black}$\checkmark$} Only modest behavioural change (+9.4\% inter-order interval, Monsoon)\\
     {\color{red!60!black}$\times$} No consistent EENS improvement (5/36 better, 11/36 worse)};

% Conclusion: full-width band beneath both columns
\node[conclusion, below=15mm of evidence, xshift=-41mm] (conclusion)
    {Conclusion\\[2pt]
     No consistent EENS improvement from EWMA.
 Existing state information may overlap with
 the delay estimate, or the added signal may
 have limited predictive value under the
 evaluated supplier model.};

% Arrows
\draw[arr]  (comment)     -- (iid);
\draw[arr]  (iid)         -- (regime);
\draw[arr]  (regime)      -- (informative);
\draw[jump] (informative.east) -- ++(6mm,0) |- (ewma.west);
\draw[arr]  (ewma)        -- (rlinvda);
\draw[arr]  (rlinvda)     -- (eval);
\draw[arr]  (eval)        -- (evidence);
\draw[arr]  (informative.south) |- (conclusion.west);
\draw[arr]  (evidence.south) |- (conclusion.east);

\end{tikzpicture}%
}
\caption{Experimental design of the RLInv-DA
delay-aware extension. A persistent two-state
supplier model introduces temporal dependence in
delivery delays. An EWMA-based delay estimate is
added to the controller observation and evaluated
against a matched base policy. The additional
observation does not produce a consistent
reliability improvement under the evaluated
supplier-regime model.}
\label{fig:rlinv_da_rationale}
\end{figure}



% ============================================================
\section[Extended Controller Comparisons]
{\revised{Extended Controller Comparisons}}
% ============================================================

\subsection[Calibrated Classical Baselines]
{\revised{Calibrated Classical Baselines}}
\label{sec:calibrated_baselines}

\begin{revisedblock}
The classical $(s,S)$ controller was evaluated using
three alternative threshold configurations:
B1-Analytic and two calibrated variants,
B1*-Deployable and B1*-HighBuffer (collectively B1*).
B1-Analytic determines its reorder point from expected
consumption over the assumed 24-hour Normal supplier
lead time. The B1* parameters were calibrated
separately for each site on days 1--45 under Normal
delivery logistics, without using the test window for
threshold selection.

The calibration examined 503 candidate parameter pairs
in total across ten sites (450 coarse-grid candidates
and 53 distinct refinements). Each candidate was scored
on ten stochastic training-window episodes, giving
5,030 candidate-evaluation episodes. Initial inventory
was sampled uniformly from 20--90\% of tank capacity.
B1*-Deployable first minimised mean training EENS,
then considered diesel use and order frequency within
a narrow near-best EENS band. B1*-HighBuffer selected
the highest reorder point among candidates achieving
the exact minimum training EENS, using the previously
recorded calibration evaluations without additional
search episodes. Neither variant changes its thresholds
during held-out operation.

Tables~\ref{tab:calibrated_baselines} and
\ref{tab:calibrated_baseline_orders} report the
held-out comparison over ten sites, ten evaluation
seed blocks and ten episodes per seed and delivery
scenario. For the deterministic classical controllers,
seed blocks represent evaluation randomness rather
than independently trained policies. The four-algorithm
comparison in Subsection~\ref{sec:g4_rl_comparison}
uses five training seeds; its RLInv means therefore
have a different aggregation population.
\end{revisedblock}

\begin{table}[htbp]
\centering
\small
\caption{\revised{Mean EENS (kWh per 360-hour episode)
for classical ordering baselines and RLInv: ten sites,
ten evaluation seed blocks, and ten episodes per block
and delivery-logistics scenario. Lower is better.}}
\label{tab:calibrated_baselines}
\begin{tabular}{lrrrr}
\toprule
\textbf{Controller} & \textbf{Normal} & \textbf{Delayed} &
\textbf{Monsoon} & \textbf{Extreme} \\
\midrule
\revised{B1-Original} & \revised{3.5476} & \revised{13.4320} &
\revised{23.2043} & \revised{57.7315} \\
\revised{B1-Analytic} & \revised{1.6862} & \revised{12.4825} &
\revised{23.1672} & \revised{61.9676} \\
\revised{B1*-Deployable} & \revised{0.6744} & \revised{7.1170} &
\revised{16.4545} & \revised{55.6580} \\
\revised{B1*-HighBuffer} & \revised{0.0628} & \revised{3.3952} &
\revised{11.2822} & \revised{53.1643} \\
\midrule
\revised{RLInv} & \revised{0.0189} & \revised{3.5992} &
\revised{11.2865} & \revised{53.8544} \\
\bottomrule
\end{tabular}
\end{table}

\begin{table}[htbp]
\centering
\small
\caption{\revised{Mean positive orders accepted per 360-hour
episode across ten sites and ten evaluation seed blocks,
under the same conditions as Table~\ref{tab:calibrated_baselines}.}}
\label{tab:calibrated_baseline_orders}
\begin{tabular}{lrrrr}
\toprule
\textbf{Controller} & \textbf{Normal} & \textbf{Delayed} &
\textbf{Monsoon} & \textbf{Extreme} \\
\midrule
\revised{B1*-Deployable} & \revised{0.982} & \revised{0.938} &
\revised{0.884} & \revised{0.637} \\
\revised{B1*-HighBuffer} & \revised{1.998} & \revised{1.954} &
\revised{1.829} & \revised{1.252} \\
\revised{RLInv} & \revised{2.164} & \revised{2.106} &
\revised{1.970} & \revised{1.341} \\
\bottomrule
\end{tabular}
\end{table}

\begin{revisedblock}
Calibration substantially changes the classical
controller's reliability. In particular,
B1*-HighBuffer records numerically lower mean EENS
than RLInv in three of the four delivery scenarios:
Delayed (3.3952 versus 3.5992), Monsoon (11.2822
versus 11.2865), and Extreme (53.1643 versus
53.8544\,kWh per episode). RLInv retains the lowest
mean under Normal logistics (0.0189 versus 0.0628).
HighBuffer also places fewer accepted orders than
RLInv in all four scenarios. Thus, on this benchmark,
the proposed learned controller does not establish a
consistent numerical reliability or ordering-frequency
advantage over a carefully calibrated classical rule.
The mean differences do not constitute a statistical
test of superiority, and order count is not a complete
operating-cost measure.

HighBuffer improves substantially on Deployable under
Delayed and Monsoon logistics, but its advantage is not
universal. On site5 in the ten-seed Extreme-scenario
classical comparison, Deployable records 241.348\,kWh
EENS versus 243.652\,kWh for HighBuffer. The selected
HighBuffer thresholds came from the training-window
candidate grid, not from the observed test results.
A limited site5 held-out pilot was inspected after the
rule was fixed and before the full evaluation; the
pilot observations were included in the subsequently
reported full grid, and the thresholds were not changed.
This disclosure prevents the pilot from being mistaken
for an independent validation sample.

The results identify threshold calibration and buffer
selection as important drivers of reliability under
uncertain diesel delivery. They also place a clear
boundary on the measured incremental advantage of
jointly learned ordering and dispatch.
\end{revisedblock}

\subsection[Comparison of Reinforcement Learning Algorithms]
{\revised{Comparison of Reinforcement Learning Algorithms}}
\label{sec:g4_rl_comparison}

\begin{revisedblock}
Across the ten-site evaluation, RLInv, unmasked PPO (A7),
Deep Q-Network (DQN), and Advantage Actor--Critic (A2C)
achieve broadly similar reliability. The most pronounced
behavioural distinction is in how their ordering actions
interact with the inventory-feasibility rules.

Each controller was trained separately for each site using
five seeds (42, 123, 777, 7, and 21), with a nominal budget
of 400,000 environment steps per checkpoint. RLInv and A7
completed 401,408 steps to accommodate whole training rollouts.
Evaluation followed the held-out protocol in
Section~\ref{sec:evaluation_protocol}: days 46--60,
four delivery-logistics scenarios, and ten episodes per
site--seed--scenario. Initial inventory was sampled from
30--90\% of tank capacity during evaluation, and each episode
lasted 360 hourly steps. The study therefore comprises
200 site--controller--seed checkpoints, 8,000 evaluation
episodes, and 2,880,000 simulated hourly steps.
\end{revisedblock}

\begin{table}[htbp]
\centering
\small
\caption{\revised{Mean EENS (kWh per 360-hour episode)
across all ten sites, five training seeds, and ten evaluation
episodes per seed and delivery scenario. Lower is better.}}
\label{tab:g4_rl_ten_site}
\begin{tabular}{lrrrr}
\toprule
\textbf{Controller} & \textbf{Normal} & \textbf{Delayed} &
\textbf{Monsoon} & \textbf{Extreme} \\
\midrule
\revised{RLInv (MaskablePPO)} & \revised{0.019} & \revised{3.119} &
\revised{12.668} & \revised{57.072} \\
\revised{A7 (PPO)} & \revised{0.027} & \revised{3.557} &
\revised{12.785} & \revised{57.454} \\
\revised{DQN} & \revised{0.019} & \revised{3.074} &
\revised{12.153} & \revised{56.789} \\
\revised{A2C} & \revised{0.122} & \revised{3.486} &
\revised{12.766} & \revised{57.082} \\
\bottomrule
\end{tabular}
\end{table}

\begin{revisedblock}
The RLInv means here use five training seeds per site,
whereas the calibrated-classical comparison in
Table~\ref{tab:calibrated_baselines} uses ten.
The respective RLInv Extreme-scenario estimates are
57.072 and 53.8544\,kWh per episode. These are
different evaluated seed populations, not discrepant
entries for an identical aggregation. The observed
five-seed mean is approximately 6\% higher; this
illustrates sensitivity to the choice of evaluation
population, without identifying a unique cause for
the difference.
\end{revisedblock}

\begin{revisedblock}
Table~\ref{tab:g4_rl_ten_site} shows that DQN has the
lowest numerical mean EENS under Delayed, Monsoon, and
Extreme logistics, while RLInv records lower means than
A7 and A2C under those scenarios. In 19 of the 40
site--scenario combinations, EENS is exactly zero for
all four controllers. These exact ties limit how much
of the evaluation grid can distinguish their reliability.

Exploratory paired comparisons use the site as the unit
of analysis: evaluation episodes and training seeds are
averaged within each site, yielding ten paired site-level
means per scenario. None of the 12 RLInv-versus-comparator
comparisons is statistically significant after Holm
correction, and all non-degenerate 95\% confidence intervals
include zero. In the Normal RLInv--DQN comparison, all ten
paired site differences are exactly zero; the paired
\emph{t}-test is undefined and its \emph{p}-value is
reported as not applicable.
\end{revisedblock}

\begin{table}[htbp]
\centering
\small
\caption{\revised{EENS on site5, the site with the highest
Extreme-scenario unmet energy in the ten-site comparison.
Values are mean kWh per 360-hour episode across five seeds
and ten evaluation episodes per seed and scenario.}}
\label{tab:g4_rl_site5}
\begin{tabular}{lrrrr}
\toprule
\textbf{Controller} & \textbf{Normal} & \textbf{Delayed} &
\textbf{Monsoon} & \textbf{Extreme} \\
\midrule
\revised{RLInv} & \revised{0.000} & \revised{12.464} &
\revised{58.509} & \revised{257.897} \\
\revised{A7} & \revised{0.000} & \revised{15.911} &
\revised{61.653} & \revised{260.711} \\
\revised{DQN} & \revised{0.000} & \revised{12.639} &
\revised{55.240} & \revised{255.662} \\
\revised{A2C} & \revised{0.722} & \revised{16.452} &
\revised{60.819} & \revised{261.194} \\
\bottomrule
\end{tabular}
\end{table}

\begin{revisedblock}
Site5 provides a demanding contrast to the many exact-zero
operating points. Under Monsoon and Extreme logistics,
DQN has numerically lower EENS than RLInv on this site,
although the aggregate paired comparisons above remain
inconclusive. Under Normal logistics, A2C has nonzero
mean EENS at four sites (site3, site5, site7, and site8).
Site5 contributes approximately 59\% and site3 25\% of
A2C's ten-site Normal mean; the difference is therefore
concentrated in a small number of sites.
\end{revisedblock}

\begin{table}[htbp]
\centering
\scriptsize
\setlength{\tabcolsep}{3pt}
\caption{\revised{Ordering and energy outcomes under Normal
logistics for the trace-audited sites site2, site4, site6,
site7, site8, site9, and site10 (350 episodes per
controller). Quantity is accepted at order placement,
not confirmed delivery. Diesel denotes energy consumed.}}
\label{tab:g4_order_quantity_subset}
\begin{tabular}{lrrrrrr}
\toprule
\textbf{Controller} &
\shortstack{\textbf{Orders}\\\textbf{per ep.}} &
\shortstack{\textbf{Mean order}\\\textbf{size (kWh)}} &
\shortstack{\textbf{Accepted qty.}\\\textbf{(kWh/ep.)}} &
\shortstack{\textbf{Mean tank}\\\textbf{inv. (\%)}} &
\shortstack{\textbf{Diesel}\\\textbf{(kWh/ep.)}} &
\shortstack{\textbf{EENS}\\\textbf{(kWh/ep.)}} \\
\midrule
\revised{RLInv} & \revised{2.114} & \revised{137.385} &
\revised{290.47} & \revised{79.4} & \revised{194.190} & \revised{0.027} \\
\revised{A7} & \revised{8.566} & \revised{38.820} &
\revised{332.52} & \revised{89.9} & \revised{194.178} & \revised{0.039} \\
\revised{DQN} & \revised{9.091} & \revised{36.727} &
\revised{333.91} & \revised{91.6} & \revised{194.190} & \revised{0.027} \\
\revised{A2C} & \revised{8.517} & \revised{39.076} &
\revised{332.82} & \revised{91.0} & \revised{194.190} & \revised{0.027} \\
\bottomrule
\end{tabular}
\end{table}

\begin{table}[htbp]
\centering
\scriptsize
\setlength{\tabcolsep}{3pt}
\caption{\revised{Order-action waterfall under Normal
logistics for the same seven sites (350 episodes per
controller). The size-reduced column is a subset of
positive accepted orders, not an additional outcome.}}
\label{tab:g4_order_waterfall}
\begin{tabular}{lrrrrr}
\toprule
\textbf{Controller} &
\shortstack{\textbf{Order}\\\textbf{requests}} &
\shortstack{\textbf{Blocked}\\\textbf{pending}} &
\shortstack{\textbf{Clipped}\\\textbf{to zero}} &
\shortstack{\textbf{Positive}\\\textbf{accepted}} &
\shortstack{\textbf{Accepted:}\\\textbf{size-reduced}} \\
\midrule
\revised{RLInv} & \revised{740} & \revised{0} & \revised{0} &
\revised{740} & \revised{0} \\
\revised{A7} & \revised{81,747} & \revised{50,586} &
\revised{28,163} & \revised{2,998} & \revised{2,804} \\
\revised{DQN} & \revised{83,219} & \revised{51,211} &
\revised{28,826} & \revised{3,182} & \revised{3,012} \\
\revised{A2C} & \revised{84,885} & \revised{54,412} &
\revised{27,492} & \revised{2,981} & \revised{2,795} \\
\bottomrule
\end{tabular}
\end{table}

\begin{revisedblock}
The trace-level ordering diagnostic in
Tables~\ref{tab:g4_order_quantity_subset} and
\ref{tab:g4_order_waterfall} explains the different
replenishment patterns. Under Normal logistics, all
four controllers consume approximately 194.2\,kWh of
diesel energy per episode and have near-identical EENS.
RLInv accepts 2.114 orders per episode with a mean size
of 137.385\,kWh, against 8.517--9.091 smaller orders
for the unmasked controllers. RLInv also holds less
inventory on average.

Of A7's 81,747 attempted order actions, 50,586 were
blocked while a delivery was pending and 28,163 were
clipped to zero because the available tank space did
not permit a positive order. The remaining 2,998
requests yielded positive accepted orders, including
2,804 whose quantity was reduced at placement. Thus
96.3\% of A7's requests yielded no positive order.

The mechanism follows from the implemented feasibility
interface. RLInv's pre-action mask permits an order only
when the full requested quantity fits in the tank and
no delivery is pending; the unmasked controllers can
request orders that the simulator subsequently blocks
or reduces. RLInv's absence of blocked or clipped
requests is therefore a property of masking, not a
separate demonstration of learned feasibility.
Accepted order quantities are distinct from confirmed
fuel deliveries and purchases.
\end{revisedblock}

\begin{revisedblock}
These results also inform the choice of MaskablePPO
in Section~\ref{sec:rationale}. In the evaluated
simulator, the physical safety layer prevents infeasible
orders from executing even when an unmasked controller
requests them. The masked configuration shows a clear
difference in ordering behaviour but no statistically
distinguishable EENS advantage over the evaluated
unmasked configurations. Thus, the demonstrated value
of pre-action masking here lies in avoiding infeasible
requests and their subsequent correction, rather than
in an observed reliability gain.
\end{revisedblock}

\begin{revisedblock}
A supplementary training-budget analysis compared
DQN and A2C checkpoints after 50,000, 100,000, and
200,000 environment steps with their final 400,000-step
checkpoints. On site5, DQN approached its final observed
EENS by 100,000 steps, whereas A2C showed non-monotonic
performance. The checkpoint-level confidence intervals
and protocol details are provided in
\hyperref[app:rl_training_budget]{Appendix~\ref*{app:rl_training_budget}}.
\end{revisedblock}

\begin{revisedblock}
The ten-site comparison was exploratory, and the
observed differences describe trained controller
configurations rather than isolated algorithmic effects.
RLInv was trained with initial inventory sampled uniformly
over 20--90\% of tank capacity, whereas A7, DQN, and A2C
used a fixed initial fraction of 60\%; inventories varied
during all training episodes. Algorithm implementations,
optimisation settings, and the feasibility interfaces also
differed. These factors particularly limit causal
interpretation of ordering-frequency and inventory results.
The original RLInv/A7 training-source contemporaneity
qualification remains despite independent replay of the
evaluation trajectories.
\end{revisedblock}

\begin{revisedblock}
Statistical non-significance does not establish
equivalence. Order counts and accepted quantities do
not independently establish lower operating expenditure,
and this four-algorithm study does not establish a general
reliability advantage over the calibrated classical
controllers in Subsection~\ref{sec:calibrated_baselines}.
\end{revisedblock}




% ============================================================
\section[Computational Cost and Cross-Site Generalisation]
{\revised{Computational Cost and Cross-Site Generalisation}}
\label{sec:computational_generalisation}
% ============================================================
Scalability is examined from two perspectives:
the computational cost of executing individual
controller decisions and the ability to reuse
a trained policy across previously unseen sites.
These evaluations do not constitute an
end-to-end multi-tower deployment benchmark.

\subsection[Computational Cost of the Controllers]
{\revised{Computational Cost of the Controllers}}
\label{sec:controller_computational_cost}

\begin{revisedblock}
In addition to reliability, the computational cost
of a controller is relevant to its practical
application. The controllers were therefore
examined in terms of decision latency, model size,
and recorded training effort.

The measurements were performed on a Windows 11
computer with a six-core AMD processor and
approximately 8 GB of RAM, using CPU execution.
For the four RL controllers, 15 trained models
per algorithm were benchmarked across site1,
site3, and site5, with five training seeds per
site. Each measurement used 100 warm-up calls
followed by 1,000 timed decisions.


Table~\ref{tab:rl_computational_cost} reports
decision latency across the 15 trained models
of each RL algorithm. For each model, 1,000
decisions were timed, and the median latency
was calculated. The reported value is the
median across these 15 model-level results.

Two measurements are reported: policy-call
latency, which times the model prediction
interface, and controller-call latency, which
additionally includes observation 
normalisation
and the implemented action-selection interface.
Both exclude simulator stepping and external
communication.


\begin{table}[htbp]
\centering
\small
\caption{\revised{Measured computational characteristics
of the four RL controllers. Decision times are
reported in milliseconds; lower latency indicates
faster execution under the measured conditions.}}
\label{tab:rl_computational_cost}
\begin{tabular}{lrrrr}
\toprule
\textbf{Controller} &
\textbf{Parameters} &
\textbf{Size (MiB)} &
\textbf{Policy (ms)} &
\textbf{Controller (ms)} \\
\midrule
\revised{RLInv} &
\revised{139,527} &
\revised{1.628} &
\revised{1.182} &
\revised{1.421} \\
\revised{A7} &
\revised{139,527} &
\revised{1.629} &
\revised{0.640} &
\revised{0.867} \\
\revised{DQN} &
\revised{140,812} &
\revised{1.102} &
\revised{0.372} &
\revised{0.538} \\
\revised{A2C} &
\revised{139,527} &
\revised{1.629} &
\revised{0.651} &
\revised{0.863} \\
\bottomrule
\end{tabular}
\end{table}
\end{revisedblock}

\begin{revisedblock}
All four RL controllers have approximately
140,000 trainable parameters, and their saved
model files occupy approximately 1.1--1.6 MiB.

Each controller's action-producing network
contains 70,406 parameters, although the
complete training architectures differ.


DQN records the lowest measured controller
latency at 0.538 ms, followed by A2C at
0.863 ms, A7 at 0.867 ms, and RLInv at
1.421 ms. These measurements demonstrate that
individual decisions can be computed within
a few milliseconds on the tested CPU.

The timings are implementation-specific.
They include different algorithm interfaces
and, for unmasked controllers, additional
handling within the evaluation interface.
The measured observations also arise from
policy-specific trajectories. Consequently,
the timing differences cannot be attributed
solely to the learning algorithm or to
action masking.
\end{revisedblock}

\begin{revisedblock}
The computational comparison was extended to
the original rule-based inventory controller
(B1), the 24-hour rolling-horizon MPC controller,
and the perfect-information Oracle-MPC reference.
Table~\ref{tab:classical_computational_cost}
reports their measured decision times.
\end{revisedblock}

\begin{table}[htbp]
\centering
\small
\caption{\revised{Measured decision latency of classical
and optimisation-based reference controllers.
Oracle-MPC was assessed using a smaller
diagnostic sample.}}
\label{tab:classical_computational_cost}
\begin{tabular}{lrr}
\toprule
\textbf{Controller} &
\textbf{Median (ms)} &
\textbf{Timed decisions} \\
\midrule
\revised{B1-Original} &
\revised{0.0181} & \revised{1,000} \\
\revised{MPC ($H=24$ h)} &
\revised{8.511} & \revised{1,000} \\
\revised{Oracle-MPC ($H=360$ h)} &
\revised{500.707} & \revised{20} \\
\bottomrule
\end{tabular}
\end{table}

\begin{revisedblock}
The original B1 threshold controller has the
lowest measured decision latency, as expected
from its simple rule-based implementation.
The 24-hour MPC controller requires more
computation because it solves an optimisation
problem at each decision. The substantially
longer Oracle-MPC horizon and its access to
future information make it a different,
offline reference problem. Its latency is
based on only 20 diagnostic calls and is
not directly comparable in measurement
precision with the other controllers.

The classical-controller timings use the
original fixed B1 thresholds, not the
site-calibrated B1* configurations examined
in Section~7.6.1. All reported measurements
exclude environment stepping and external
communication overhead. They therefore
characterise local decision computation,
rather than complete deployment latency.
\end{revisedblock}

\begin{revisedblock}
Recorded production training times also show
differences in computational effort. The median
training duration was 53.69 minutes for DQN
and 5.16 minutes for A2C under the recorded
training configuration. These values describe
the implemented training runs, not an
algorithm-independent speed comparison.

The historical RLInv timing record includes
an MLflow session interval rather than an
equivalent production training-duration
measurement, while the original A7 training
duration could not be independently established.
They are therefore excluded from the direct
training-time comparison.

Overall, the benchmarks show that the evaluated
RL controllers require only millisecond-scale
computation per decision on the tested CPU.
This supports the computational practicality
of the proposed tower-level control approach.

The measurements are limited to local controller
execution and do not include data acquisition,
external communication, or physical deployment.
Cross-site generalisation of the learned
control policy is examined separately in the
following subsection.

\end{revisedblock}


% ============================================================
\section{Summary}
\label{sec:results_summary}
% ============================================================
\begin{sloppypar}
Within the original policy and ablation comparisons,
RLInv records the lowest mean EENS under Normal and
Extreme logistics, is statistically indistinguishable
from the closest comparator under Delayed logistics,
and is numerically outperformed by A6 under Monsoon
logistics.

However, the expanded calibrated-classical comparison
shows that B1*-HighBuffer achieves lower mean EENS
than RLInv under Delayed, Monsoon, and Extreme logistics,
while RLInv retains the lower mean under Normal logistics.
B1*-HighBuffer also places fewer orders than RLInv
across all four scenarios.

These findings demonstrate that careful calibration
of classical inventory-ordering thresholds can match
or outperform the proposed reinforcement learning
controller under stressed delivery conditions. The ablations provide a
more qualified explanation than the primary comparison alone. The
additional benefit of jointly learned replenishment is
scenario-dependent, appearing under Normal and Extreme conditions but
not under Delayed or Monsoon logistics. The cleaner RLInv--A7
comparison finds no statistically distinguishable mean-EENS effect
from action masking when replenishment is jointly learned, while the
larger A6--TrackB gap under stressed logistics is confounded by
differences in observation design. Removing the five logistics-related
observation channels produces only small, non-significant changes.
Reward, tank-capacity, and delivery-distribution analyses further
clarify the conditions under which RLInv's advantage persists, while
the stochastic-grid experiment and RLInv-DA extension identify
limitations rather than universal improvements.
\end{sloppypar}